% Related work: position the envelope among the literatures it composes.
The framework sits downstream of most machinery that cross-domain scientific
integration needs. Extraction, schema matching, entity resolution, fusion,
contradiction and stance analysis, and literature-based discovery are treated
as available capability. It also imports partial identification and robust
decision theory rather than claiming either. The residual question is their
composition: after heterogeneous sources have been structured and aligned,
which global claims are invariant across all source-compatible completions,
which differences certify obstruction or plurality, and what downstream action
is justified when scientific identity remains partially identified?

\subsection{Partial identification, imprecise probability, and decisions}

Partial-identification theory replaces an unjustified point estimate with the
set of parameter values compatible with observations and assumptions
\cite{manski2003partial,tamer2010partial}. Imprecise-probability and robust
Bayesian traditions similarly represent a set of probability models and
choose actions against the resulting range of expected losses
\cite{walley1991imprecise,berger1985decision}. Blackwell's comparison of
experiments orders information structures by the decisions they support
\cite{blackwell1953equivalent}. These are direct conceptual donors. The
epistemic portrait envelope transports their logic from parameters and
experiments to provenance-bound, typed scientific representations: the
identified object is a set of global portraits or claim values; an observation
refinement shrinks that set; and a downstream loss is declared separately from
the scientific relation being represented.

The distinction matters because abstention is not automatically a virtue. It
is optimal only under a loss model and an identified set that make deferral
cheaper than the worst supported merge or split. Conversely, a determinate
mapping is not licensed merely because a model is confident. Point
identification is a property of the observation fibre and assumptions, not of
the sharpness of a model's posterior. The new contribution is thus not
set-valued inference in general, but its composition with scientific
coordinate preservation, obstruction, provenance and global portrait
construction.

\subsection{Constraint composition and global consistency}

Ontology and schema matchers commonly propose pairwise correspondences, while
constraint-processing research distinguishes local consistency from the
existence of a joint solution~\cite{dechter2003constraint}. The envelope uses
that distinction as an authority boundary. A candidate mapping contributes a
constraint on feasible worlds; it does not authorize a global identity until
all accepted constraints possess a common completion and the claim is stable
over those completions. This absorbs generic constraint satisfaction as a
donor mechanism while giving its global-consistency result a scientific-
identity interpretation.

\subsection{Structured scientific knowledge bases and extraction}

MUSE~\cite{muse2026} constructs a full-text, multi-domain resource of
source-grounded problem--solution--rationale structures, and the Discovery
Engine~\cite{discoveryEngine2025} distils literature into structured knowledge
artifacts and an integrated computational landscape. Both target the
representation itself: what a source contains, recovered at scale. The
decision studied here begins after that representation exists, and asks
whether referent, construct, measurement, context, modality and attribution
distinctions can govern a downstream mapping outcome and produce an
obstruction when compatibility has not been earned. SciER~\cite{scier2024}
supplies scientific entity and relation annotations for datasets, methods and
tasks in full text, which makes information extraction, discourse processing
and entity/relation recovery an input to this work rather than part of it.

\subsection{Schema-centric approaches}

SciSchema.org~\cite{scischema2026} provides multidisciplinary,
expert-reviewed scientific-process schemas with parameters, measurements,
procedures and provenance, and is strong prior art for reusable scientific
schema design. SCOPE and SCION~\cite{scopes2026} introduce an auditable
benchmark and reference pipeline for evidence-linked schema induction and
optional fusion from text, including conservative fusion behaviour. Executable
Schema Contracts~\cite{schemaContracts2026} discover an executable schema from
heterogeneous sources and use it to drive provenance-aware construction and
multi-source retrieval. Each of these produces a shared structure that a
downstream consumer can rely on; none of them is asked whether a particular
proposed mapping preserves the coordinates a specific scientific claim depends
on. That is the discriminator evaluated here, over already-structured
projections and including cases in which refusing a merge is the correct
outcome.

\subsection{Alignment, integration, and resolution}

Ontology matching~\cite{euzenat2007ontology} studies correspondences such as
equivalence, subsumption and disjointness between ontology entities, and
language-model systems such as LLMATCH~\cite{llmatch2025} add modern
schema-preparation, candidate-selection and column-alignment machinery.
Surveys of knowledge-graph construction~\cite{hofer2024kgconstruction} cover
heterogeneous-source integration, ontology and schema matching, entity
resolution, fusion, provenance and quality assurance, while reviews of
scholarly knowledge graphs~\cite{verma2023scholarlykg} cover extraction,
construction, refinement and use of semantic scholarly infrastructures. The
distinction pursued here is that a mapping decision is bound to several
scientific coordinates at once, and that the exact source, provenance and
non-preserved-coordinate information is retained with the resulting decision.
Entity resolution and fusion are especially useful pressure cases, because a
record-level match may be appropriate at exactly the moment a claim-level
merge remains unsafe on construct, measurement, context or modality grounds.

\subsection{Claim comparison and stance}

Stance detection assigns target-relative labels such as support or opposition;
surveys document its use in fact-checking and misinformation
settings~\cite{hardalov2022stance}. Apparent scientific disagreement, however,
frequently arises from different context, modality, measurement or attribution
coordinates rather than from a genuine opposition, and a mapping record that
collapses every difference into one generic contradiction label discards the
reason. Preserving that reason, rather than detecting stance, is what is
evaluated here.

\subsection{Scientific-variable semantics and interoperability}

The I-ADOPT interoperability framework~\cite{iadopt2021} provides a
community-developed semantic decomposition for machine-readable observable and
scientific-variable descriptions, and a recent benchmark builds an
expert-annotated corpus of more than one hundred scientific variables and
reports that large language models still struggle to construct those
structured representations accurately~\cite{iadoptbench2026}. FAIR~2.0
separates terminological, ontological, referential, schema and logical
interoperability, and argues that semantic interoperability requires explicit
mappings rather than one privileged ontology~\cite{fair2_2024}. Recent
provenance work goes further: DEC reads provenance-bearing statements as
epistemic stance and preserves source-relative disagreement instead of
collapsing every attributed assertion into a single factual
core~\cite{dec2026}. Structured variable semantics, generic semantic
interoperability and provenance-aware pluralism are therefore supplied by this
literature. The residual question is claim-relative and downstream of all of
it: once a strong semantic integration stack has produced structured candidate
mappings, what further conditions license an assertion of \emph{scientific
identity}? The decision contract studied here consumes ontology and schema
alignment, provenance, plural-view machinery and consistency checks as inputs,
and returns a glue, obstruction, plural-view or unresolved outcome from
load-bearing construct, measurement, context, provenance-stance and
global-composition conditions.

\subsection{Measurement equivalence and construct validity}

Measurement-equivalence and replication literatures established long ago that
nominally identical outcome measures need not be comparable across settings,
and that equivalence can require explicit testing~\cite{jilke2017measurement}.
The mechanism here makes measurement and operationalization one explicit
coordinate of the mapping record and allows an unresolved measurement relation
to block a global merge on its own.

\subsection{Literature-based discovery and pluralism}

Swanson's literature-based discovery
lineage~\cite{swanson1986,swanson1997interactive} shows how complementary
literatures can expose implicit connections without direct citation links. The
bounded requirement added here is on the bridge concept itself: it must retain
compatible scientific meaning across the two source-local projections, or the
bridge stays an obstruction rather than being promoted to a scientific
integration. Scientific pluralism has similarly argued that multiple
representations may remain legitimate without collapsing into one complete
account~\cite{kellert2006,cartwright1999}; the computational question is
whether an integration record can carry a plural or obstructed outcome, with
exact source lineage, instead of forcing a canonical merged statement.

\subsection{Scientific retrieval and cross-domain synthesis}

OpenScholar~\cite{openScholar2024} performs retrieval-augmented synthesis
over the scientific literature at scale, and BioSage~\cite{biosage2025} combines
retrieval, specialized agents, cross-disciplinary translation and reasoning
for scientific tasks. Retrieval, synthesis, agent orchestration and
cross-domain translation are baseline capability for the problem, as are
language-model-assisted schema matching and induction. The experiment reported
here deliberately holds all of that fixed and varies only the decision rule
applied after scientific projections are already structured, so that the
construct-validity question can be asked without a retrieval confound.

The general claim is theoretical rather than an empirical promotion: once the
donor mechanisms leave more than one scientifically admissible completion, a
single canonical portrait requires either additional information, an explicit
assumption, or a declared decision rule. The current external evidence is much
narrower. On a disjoint, prospectively frozen public-reference holdout of
already-structured projections, the typed mapping calculus reduced false
merges against flat predicate canonicalization while satisfying the
false-split guard declared in advance. Raw-text extraction, naturalistic
envelope validity, independent expert construct validity, recoverability and
downstream utility require separate evidence.
\paragraph{Uncertainty in automated ontology matching.}
Osman, Pileggi and Ben Yahia empirically document substantial uncertainty and
unreliability in automated ontology matching on a real-data integration case,
and argue that appropriately designed semi-supervised approaches are more
mature than unqualified automation \cite{osman2024uncertainty}. We absorb that
diagnosis rather than claim uncertainty awareness as novel. Their unit is
matcher uncertainty and its management. Our residual is claim-relative
identification: content uncertainty and missing source, rights, executable,
reference or custody authority are represented in one terminal-preserving
fibre calculus, and a missing comparison is not converted into an observed
error, tie or obstruction.

\endinput
